% ============================================================
%  Professional Statement on the Use of Artificial Intelligence
%  Ronald J. Tackett, Ph.D.
%  Engine: XeLaTeX
%
%  CHANGELOG
%  v1.6  2026-09-04  SoFS start corrected from June to July 2025.
%  v1.5  2026-09-03  Replaced the drafted verification sentence in II with
%                    Ron's actual practice: all AI-assisted code is thoroughly
%                    tested, and the prose/code distinction is made explicit.
%  v1.4  2026-09-03  Added interactive-visualization practice to II (HTML/CSS/JS
%                    built with Claude), with the two reservations and the
%                    limiting-case verification check. Tied it into VIII.
%  v1.3  2026-09-03  Added grading practice to II (names redacted, grades
%                    assigned by hand, AI used as a consistency and
%                    completeness check on already-evaluated work).
%  v1.2  2026-09-03  Clarified II and III: quantitative work is code-mediated
%                    (Claude writes Python/pandas parsers that run offline),
%                    so institutional data is not sent to the model at all.
%  v1.1  2026-09-03  Added Section III (Where I Do Not Use It): peer-review
%                    exclusion and de-identification practice. Added output
%                    verification to II. Sharpened student-access claim in IV.
%                    Updated VIII for the CoLLT gatherings. Rewrote IX.
%                    Sections renumbered III-VIII to IV-IX.
%  v1.0  2026-09-03  Rebuilt in LaTeX from the March 2026 PDF.
%                    Content unchanged; footer date moved to
%                    September 3, 2026 and the closing block from
%                    Spring 2026 to Fall 2026.
% ============================================================
\documentclass[10pt]{article}
\usepackage{tackettdoc}

\setdoclabel{AI Personal Statement}
\setupdated{September 3, 2026}

\begin{document}

\statementheader
  {Founding Head, School of Foundational Studies}
  {Associate Professor of Physics \textperiodcentered\ Kettering University}
  {Professional Statement on the Use of Artificial Intelligence}

\lede{This statement reflects where I am right now in my thinking about artificial intelligence
in my professional work. It is not a final position. I am still working through many of these
questions, and I expect this statement to look different in a year or two.}

\stsection{I. Context}

I became Founding Head of the School of Foundational Studies (SoFS) in July 2025, which means I
am standing up a new academic unit at a moment when questions about AI in higher education are
particularly unsettled. That timing shapes how I think about this. The decisions made now about
how the SoFS approaches AI will have some staying power, so I am trying to be deliberate rather
than reactive. This statement is part of that effort.

\stsection{II. How I Currently Use AI}

I use AI, primarily Claude (Anthropic's AI assistant), regularly across most areas of my work.
On the administrative side, this includes producing documents like budget reports, curriculum
memos, faculty communications, and enrollment analyses. Much of the quantitative work is
indirect, and deliberately so. Rather than hand institutional data to an AI system, I use
Claude to help me write Python programs, usually built on pandas, that parse the standard
reports I download from our databases and generate formatted output. Those programs run offline
on my own machine, so the data never leaves it. What I am asking AI for in that arrangement is
the tooling, not access to the records. In communications, I use AI
to refine drafts of memos, letters, and newsletter columns; I bring the substance, and AI helps
shape them. For curriculum work, I have used it to help organize reports and recommendations
for the SoFS Curriculum Council and the University Curriculum Committee.

I also use AI in building course materials and syllabi, and in the early stages of proposal
work, where it helps me outline an argument and pressure-test a draft before it goes anywhere.
The throughline in all of this is that AI handles execution on tasks where I have already done
the thinking. The judgment, the context, and the decisions remain mine.

Grading deserves its own account, because it is the use students have the most legitimate
interest in knowing about. Student names are redacted before any work goes into a tool. I assign
every grade myself. What I use AI for is a second pass over work I have already evaluated:
checking that I applied the same standard to the twentieth submission as to the first, and that
I did not overlook something a student actually did. Consistency across a stack of papers is
precisely where human attention degrades, and a second reader helps. The judgment stays mine.
What I am checking is whether I exercised it evenly.

The newest use, and the one students actually see, is building interactive visualizations of
scientific concepts. Working with Claude in HTML, CSS, and JavaScript, I can put together a
functioning simulation far faster than I could on my own: a field that reconfigures as a student
varies a parameter, a distribution that shifts as a constraint changes. These are not
demonstrations that do the reasoning for students. They make a concept confrontable while
leaving the work intact, because the student still has to predict what will happen and then find
out whether they were right. Of everything described here, this is the use I would point to as
AI serving learning rather than merely not damaging it.

Two things about it give me pause. The first is that what limits me is no longer my facility
with JavaScript; it is my judgment about whether students need the thing at all. Building has
become cheap enough that I have to watch for building because I can rather than because a
concept called for it. The second is that I am putting code in front of a class that I did not
write line by line, and a visualization carrying a subtle physics error is worse than no
visualization, because it teaches the error with the authority of a picture. Prose I spot-check, as I
said above. Code is different, because a bug does not announce itself the way an awkward
sentence does. So I thoroughly test every program I write with AI assistance, the
visualizations and the parsing scripts alike, before it does any real work. Code I did not write
line by line gets more scrutiny from me, not less.

What I do not do is take the output on faith. I spot-check rather than verify every line, and I
treat whatever goes out under my name as mine regardless of how it was produced. ``The tool
generated it'' is not a defense I am willing to offer, and building a practice that would
require me to offer it seems like a bad trade.

\stsection{III. Where I Do Not Use It}

There are places I keep these tools out entirely, and I think the exclusions say as much about
a practice as the uses do.

The first is peer review. I have served on an NSF review panel and I review for several
journals, and I do not put proposal or manuscript content into an AI system. NSF prohibits it
outright. The reasoning behind that prohibition, that material submitted in confidence should
not be fed into a system outside the review process, holds regardless of whether a rule compels
it. My own proposals are a different matter entirely; there I use AI to help outline an argument
and strengthen a draft, because that work is mine to share.

The second is anything that identifies a person. Some of the administrative work described
above touches personnel matters, student records, and student work. Much of the time the question does not
arise, because the work is done by code running locally and the records never go anywhere. When
I do need to reason through the content of something directly, it is de-identified or aggregated
first, so what I am actually discussing is a pattern rather than a person. This is a real
constraint on what AI can do for me, and there are analyses I simply do
not run as a result. I would rather accept the constraint than leave confidentiality to my
judgment in the moment, when the convenient answer and the right one are hard to tell apart.

\stsection{IV. The Benefits I See}

The most practical benefit is capacity. My role is broad, and AI allows me to produce work at a
level of quality and volume that would otherwise be difficult to sustain. For colleagues, I
think the value is similar: less friction on editing and administrative tasks, more time for the
work that actually requires human presence. For students, the benefit is real but narrower than it
is usually described. A student who has already worked a problem and wants a second explanation
at eleven at night is genuinely well served. A student who has not yet attempted the problem is
not; the same tool that would have helped an hour later will, at that moment, quietly do the work
that was supposed to build the capability. The value is not in the access itself but in where the
student is standing when they reach for it.

\stsection{V. My Concerns}

My main concern is that AI makes it easy to skip the productive struggle that builds real
competence. In foundational education, especially, the difficulty is not incidental; it is the
mechanism (it's the process, not the product, so to speak). Students who reach for AI before
they have developed enough judgment to evaluate the output are not being helped; they are being
shortchanged. I also worry about institutional policy drifting so far behind actual practice
that it becomes meaningless, and about colleagues being left to figure all of this out without
much support.

\stsection{VI. The Principles I Try to Follow}

A few things guide how I use AI in practice.

\begin{itemize}[leftmargin=1.4em, itemsep=0.35em, topsep=0.3em, parsep=0pt]
  \item I try to be transparent about when and how I use it. I do not think there is anything to
        hide, and modeling openness seems important given my role.
  \item I keep consequential judgments human. Decisions that affect people, including faculty,
        students, and the direction of the school, are mine to make, not to delegate.
  \item I think about fitness for purpose. Not every task benefits from AI involvement, and some
        are better done without it.
  \item I try to hold myself to the same standard I hold my students to.
\end{itemize}

I lead a school built on the value of foundational learning. I should not quietly let my own
tools erode the habits of mind I am asking students to develop.

\stsection{VII. AI in My Courses}

My AI policy varies by course and assignment, but the underlying message to students is
consistent: the difficulty is not the obstacle, it is the point. In foundational education, the
struggle to work through a problem, construct an argument, or derive a result is where learning
actually happens. A student who offloads that struggle to AI may produce a correct answer, but
they have skipped the process that builds the competence the answer is supposed to represent. I
try to make that distinction explicit with my students rather than simply restricting or
permitting AI use by rule. The policy question I keep coming back to is not whether AI is
allowed, but whether its use in a given moment serves the student's development or substitutes
for it.

\stsection{VIII. Where I Expect This to Go}

In administration, I expect to consolidate what I have built and make it accessible to others
rather than continue adding to it. In teaching, the more significant shift will come, in my
opinion, as AI-assisted problem-solving becomes part of the subject matter in some of my
courses rather than something to accommodate around the edges. The visualizations described
above are an early version of that shift. At the moment I build them and students use them, and
I am not convinced that is where it should end up.

At the school level, the work has already moved from designing a framework to using it, and it
has moved further outward than I anticipated when I first drafted this statement. In April 2026
I sat down with seven other liberal arts units at technological universities to compare how each
of us was structured and what each of us thought we were for. In September I argued to that same
group, by then considerably larger, that the question worth asking is not how to accommodate AI
but what we are actually educating students for. Making that case to peers clarified my own
thinking more than any amount of internal drafting had. It also means I am now committed in
public to positions I am still testing, which is uncomfortable in a way I have come to think is
useful.

I expect my use of AI to become more deliberate over time, though not necessarily more
extensive.

\stsection{IX. My Commitments}

I will keep saying plainly how I use these tools, including in rooms where that is not the
expected answer. I will keep the consequential judgments (about people, about curriculum, about
the direction of the school) on my side of the line, and I will keep the exclusions in Section
III intact even when they cost me something. And I will keep checking whether the habits I ask
students to build are habits my own practice actually reflects, because a school built on
foundational learning cannot be led by someone quietly exempting himself from it.

If this statement reads differently a year from now, that will be because the thinking moved. I
would rather revise it in public than leave it standing after I have stopped believing it.

\signatureblock{Fall 2026}

\end{document}
